% ===================================================================
%  TAULA N° 6 — L'Ostal per dedins. — Los Mòbles. — Lo Manjar.
%                L'Enlumenatge.
%  Traduction 2026 d'apres texto/io, controlee sur texto/fr et
%  texto/ca. Consigne : voir texto/oc/10-taula-01.tex.
%
%  Deux notes, toutes deux marquees « (*) », et deux appels : celui de
%  « babo » au bloc c1-12-1, celui de « lambrequino » au bloc c2-01-1.
%
%  LA FEMME DE CHAMBRE PREND SON (16), contre le fac-simile ido :
%  ecart declare dans outils/renvoji.py.
% ===================================================================

%%K t06-apar-1 apar
\VUsaut{6.0mm}
\VUtitre{15.0pt}{60}{TAULA N° 6}
\VUsaut{1.4mm}
\VUtitre{11.4pt}{0}{L'Ostal per dedins, los Mòbles, lo Manjar,}
\VUsaut{0.4mm}
\VUtitre{11.4pt}{0}{l'Enlumenatge.}
\VUsaut{2.2mm}

%%K t06-apar-2 apar
L'ostal es acabat. L'oncle de Joan s'es installat dins son novèl estatge, que ne
coneissèm la distribucion. Vejam ara cossí l'a moblat. Visitarem d'en primièr:

%%K t06-tit-1 sub
\VUpk{10.2pt}{40}{La Cambra.}
\VUsaut{3.0mm}

%%K t06-c1-01-1 p
1. --- Arribam a marrida ora, que la serviciala es ocupada a netejar la
\VUgras{cambra}.

%%K t06-c1-02-1 p
2. --- Sus la \VUgras{taula de toaleta} \textsuperscript{(1)} son escampilhats los
diferents objèctes amb los quals òm se lava, se penchena, en un mot s'apresta. Son
la \VUgras{cuveta} \textsuperscript{(2)} e lo \VUgras{pichièr}
\textsuperscript{(3)}, la \VUgras{bròssa de pels} \textsuperscript{(4)}, lo
\VUgras{penche} \textsuperscript{(5)}, lo \VUgras{sabon} \textsuperscript{(6)} e
l'\VUgras{esponga} \textsuperscript{(7)}; enfin, un \VUgras{flascon}
\textsuperscript{(8)} d'elixir dentifrici.

%%K t06-c1-03-1 p
3. --- Per tèrra, davant la taula de toaleta, vaicí lo \VUgras{selh}
\textsuperscript{(9)} per recampar las aigas brutas.

%%K t06-c1-04-1 p
4. --- Dins l'\VUgras{antecambra} \textsuperscript{(10)} vesèm qualques
\VUgras{penjadors} \textsuperscript{(13)} e dos \VUgras{parapluèjas}
\textsuperscript{(11)} dins lo \VUgras{portaparapluèjas} \textsuperscript{(12)}.

%%K t06-c1-05-1 p
5. --- De l'autre costat de la pòrta i a l'\VUgras{armari de miralh}
\textsuperscript{(14)}, que garda la \VUgras{linja} \textsuperscript{(15)}.

%%K t06-c1-06-1 p
6. --- Profitem que la \VUgras{cambrièira} \textsuperscript{(16)} fa lo
\VUgras{lièch} \textsuperscript{(17)} per n'estudiar las diferentas partidas. Lo
\VUgras{cadre del lièch} \textsuperscript{(20)} pòrta un bon \VUgras{somièr}
\textsuperscript{(21)} sul qual Justina pausarà lèu un \VUgras{matalàs}
\textsuperscript{(22)} de lana. Puèi prendrà de las cadièiras los dos
\VUgras{linçòls} \textsuperscript{(23)} e la \VUgras{cobèrta}
\textsuperscript{(24)}. Aprèp pausarà al cabeç del lièch lo \VUgras{traversièr}
\textsuperscript{(25)} e lo \VUgras{coissin} \textsuperscript{(26)}, que son amb
l'\VUgras{edredon} \textsuperscript{(27)} sul \VUgras{tapís de lièch}
\textsuperscript{(32)}. Un plumalh per espolsar los \VUgras{ridèus}
\textsuperscript{(19)} e lo \VUgras{cèl de lièch} \textsuperscript{(18)}, e vaquí
una partida de la feina ja facha.

%%K t06-c1-07-1 p
7. --- Puèi li caldrà arrengar un pauc la \VUgras{taula de nuèch}
\textsuperscript{(28)}, remontar lo \VUgras{revelh} \textsuperscript{(29)},
preparar lo \VUgras{lumenon} \textsuperscript{(30)} e emportar la
\VUgras{candela} \textsuperscript{(31)} per netejar lo candelièr.

%%K t06-tit-2 sub
\VUsaut{5.0mm}
\VUpk{10.2pt}{40}{Lo panièr de cordura e los espleches de la chiminèa.}
\VUsaut{3.0mm}

%%K t06-c1-08-1 p
8. --- Lo \VUgras{panièr de cordura} \textsuperscript{(34)} qu'es contra lo
\VUgras{tamboret} \textsuperscript{(33)} a tanben fòrça besonh d'èsser arrengat.
Caldrà plegar la \VUgras{brodadura} \textsuperscript{(35)}, sarrar dins la
\VUgras{taula de labor} \textsuperscript{(38)} lo \VUgras{didal}, lo
\VUgras{fial} \textsuperscript{(36)}, las \VUgras{agulhas}
\textsuperscript{(37)} e las \VUgras{cisèlas} \textsuperscript{(39)}, e barrar amb
la \VUgras{clau} \textsuperscript{(40)} lo cobèrcle de la taula.

%%K t06-c1-09-1 p
9. --- Sul \VUgras{guerridon} \textsuperscript{(41)} del mitan, Justina renovelarà
l'aiga de la \VUgras{carafa} \textsuperscript{(42)}. Pausarà de \VUgras{sucre} dins
lo \VUgras{sucrièr} \textsuperscript{(43)}, netejarà las \VUgras{pinças del sucre}
\textsuperscript{(44)} e emportarà la \VUgras{bròssa de rauba}
\textsuperscript{(46)}.

%%K t06-c1-10-1 p
10. --- Lo costat drech de la cambra li donarà mens de feina, que la
\VUgras{chèsa-longa} \textsuperscript{(47)} e son \VUgras{coissin}
\textsuperscript{(48)} son a lor plaça e que, coma fa pas freg, res s'es pas
bolegat entre los espleches de la \VUgras{chiminèa} \textsuperscript{(49)}.
\VUgras{Pala} \textsuperscript{(50)}, \VUgras{tenalhas} \textsuperscript{(51)},
\VUgras{capfuòcs} \textsuperscript{(55)} e \VUgras{gardacendres}
\textsuperscript{(56)} son nets; lo \VUgras{parafuòc} \textsuperscript{(54)} s'es
pas bolegat e la \VUgras{caissa de la lenha} \textsuperscript{(52)} es plan
provesida de \VUgras{lenha} \textsuperscript{(53)}.

%%K t06-noto-1 noto
\VUnotes{143.2mm}{%
(*) Babo = enfant pichon; anglés \textit{baby}, francés \textit{bébé}, italian
\textit{bambino}.%
}

%%K t06-noto-2 noto
\VUnotes{143.2mm}{%
(*) Lambrequino = francés \textit{lambrequin}, espanhòl \textit{lambrequines},
portugués \textit{lambrequins}, e quitament alemand e anglés
\textit{lambrequins}; valent a dire alemand, anglés, francés, portugués e
espanhòl parièr.%
}

%%K t06-c1-11-1 p
11. --- Li caldrà en mai levar la polsa de la \VUgras{pendula}
\textsuperscript{(57)}, dels \VUgras{candelabres} \textsuperscript{(58)} e dels
\VUgras{voidapòchas} \textsuperscript{(59)}, e puèi netejar lo \VUgras{miralh}
\textsuperscript{(60)}.

%%K t06-c1-12-1 p
12. --- Quant al \VUgras{brèç} \textsuperscript{(61)} de l'enfant (del babo (*)),
es ja prèst.

%%K t06-tit-3 sub
\VUsaut{5.0mm}
\VUpk{10.2pt}{40}{Lo Salon.}
\VUsaut{3.0mm}

%%K t06-c2-01-1 p
1. --- Vaicí ara lo salon. Es una \VUgras{cambra} fòrça polida. La chiminèa,
qu'es dins lo canton de drecha, es ornada d'una fina \VUgras{estatueta}
\textsuperscript{(1)} e de dos bèles \VUgras{vasses} \textsuperscript{(3)}, e
drapada d'un \VUgras{lambrequin} \textsuperscript{(2)} de velós (*).

%%K t06-c2-02-1 p
2. --- Per evitar que las belugas del fogal alucan lo tapís, s'es pausat davant
un \VUgras{gardafuòc} \textsuperscript{(4)} de coire.

%%K t06-c2-03-1 p
3. --- Al costat, un elegant \VUgras{escritòri} \textsuperscript{(5)} de dòna es
dobèrt. Aquí es que la \VUgras{dòna de l'ostal} \textsuperscript{(23)} escriu sas
letras.

%%K t06-c2-04-1 p
4. --- La fenèstra es garnida de \VUgras{ridèus prims} \textsuperscript{(6)}
recampats en \VUgras{embrassas} \textsuperscript{(8)} pendoladas als
\VUgras{cròcs} \textsuperscript{(7)}, e de grands ridèus de tela acabats per una
\VUgras{sanefa} \textsuperscript{(9)}.

%%K t06-c2-05-1 p
5. --- Dins l'embrasadura de la fenèstra, un \VUgras{cachapòt}
\textsuperscript{(11)} conten una \VUgras{planta} \textsuperscript{(10)} verda, una
palmièra magnifica que sas fuèlhas arriban gaireben fins a la \VUgras{lampa de
petròli} \textsuperscript{(12)}.

%%K t06-c2-06-1 p
6. --- L'\VUgras{abatjorn} \textsuperscript{(13)} de la lampa l'a fach Maria, l'encantadora
\VUgras{jova} \textsuperscript{(14)} qu'es al \VUgras{pianò}
\textsuperscript{(15)}. Assetada plan drecha sul \VUgras{tamboret}
\textsuperscript{(16)}, estúdia una pèça. Es fòrça abila e fòrça curosa, coma o
pròva son \VUgras{mòble de la musica} \textsuperscript{(17)}, qu'es perfièchament
arrengat.

%%K t06-tit-4 sub
\VUsaut{5.0mm}
\VUpk{10.2pt}{40}{Lo Bolegadís.}
\VUsaut{3.0mm}

%%K t06-c2-07-1 p
7. --- Son fraire Pau cèrca un \VUgras{libre} \textsuperscript{(19)} dins
l'\VUgras{estagièra} \textsuperscript{(18)}. Es un \VUgras{dròlle}
\textsuperscript{(20)} bolegadís e insuportable. En se pendolar a l'estagièra per
atenher lo libre qu'es tròp naut, risca de far tombar lo \VUgras{bust}
\textsuperscript{(21)} qu'es al dessús.

%%K t06-c2-08-1 p
8. --- Profita, per far aquela bestiesa, que sa maire es ocupada a charrar amb
una amiga, una \VUgras{dòna} \textsuperscript{(22)} qu'es venguda passar qualques
jorns amb ela. La maire es assetada sul \VUgras{canapè} \textsuperscript{(24)},
contra lo \VUgras{fautuèlh} \textsuperscript{(25)}, e son amiga es sus la
\VUgras{cadièira} \textsuperscript{(27)}, que ne vesèm pas que lo
\VUgras{respalhièr} \textsuperscript{(26)}.

%%K t06-c2-09-1 p
9. --- Lo grand \VUgras{cadre} \textsuperscript{(30)} pendolat a la paret conten lo
\VUgras{retrach} \textsuperscript{(29)} d'una parenta. Dins l'\VUgras{album}
\textsuperscript{(32)} qu'es sus la taula son amassadas las \VUgras{fotografias}
\textsuperscript{(33)} dels autres membres de la familha. Los enfants l'an
fulhejat i a un moment, que las \VUgras{cadièiras} \textsuperscript{(31)}
demòran encara acampadas a l'entorn de la taula.

%%K t06-c2-10-1 p
10. --- Lo \VUgras{vas chinés} \textsuperscript{(34)} de porcelana qu'es contra lo
\VUgras{dobrelètras} \textsuperscript{(35)} lo donèron a Maria per sa fèsta, e lo
\VUgras{paravent} \textsuperscript{(36)} qu'orna lo canton de la cambra lo portèt
de China son oncle.

%%K t06-c2-11-1 p
11. --- Alegrat pels bèles \VUgras{tablèus} \textsuperscript{(37)} pendolats a la
\VUgras{pareta} \textsuperscript{(42)}, esclairat pel \VUgras{lum de plafon}
\textsuperscript{(38)}, que sas \VUgras{ampolas} \textsuperscript{(39)} donan al
vèspre una lutz tan alègra, aqueste salon sembla fòrça agradiu. Lo mofle
\VUgras{tapís} \textsuperscript{(40)} espandit sul parquet e la \VUgras{sonalha
electrica} \textsuperscript{(41)}, tan practica, acaban de lo far fòrça comòde.

%%K t06-tit-5 sub
\VUsaut{3.6mm}
\VUpk{10.2pt}{40}{La Sala de manjar.}
\VUsaut{3.0mm}

%%K t06-c3-01-1 p
1. --- L'ora del sopar a sonat. Tota la familha, levat Maria, es acampada a
l'entorn de la taula. La sala de manjar es agradivament escalfada per una
excellenta \VUgras{estufa} \textsuperscript{(1)} de faiença, amb un \VUgras{tuble}
\textsuperscript{(2)} prim.

%%K t06-c3-02-1 p
2. --- Aquesta sala a pas gaire de mòbles. Al long de la paret pendola, dessús lo
\VUgras{banquet} \textsuperscript{(4)}, una \VUgras{fontaneta}
\textsuperscript{(3)} decorativa. Fòrça prèp i a lo \VUgras{bufet}
\textsuperscript{(5)}, sul qual se daissan los plats freds, las postres e los
diferents espleches de taula que servisson pas sulcòp.

%%K t06-c3-03-1 p
3. --- Atal vesèm, sus l'estagièra d'amont, un \VUgras{polet rostit}
\textsuperscript{(7)}, un \VUgras{peis} \textsuperscript{(8)} e una \VUgras{copa}
\textsuperscript{(10)} amb de \VUgras{frucha} \textsuperscript{(9)}. Al dejós,
vaicí lo \VUgras{fromatge} \textsuperscript{(11)}, lo \VUgras{pan}
\textsuperscript{(12)} e las \VUgras{olièiras} \textsuperscript{(13)}, que pòrtan
tot çò que cal per assasonar l'ensalada: l'òli, lo vinagre, la \VUgras{sal}
\textsuperscript{(14)} e lo \VUgras{pebre} \textsuperscript{(15)}. Enfin, sus
l'estagièra d'aval i a un \VUgras{gastèl} \textsuperscript{(17)} contra lo
\VUgras{mostardièr} \textsuperscript{(16)}.

%%K t06-c3-04-1 p
4. --- Lo relòtge de la sala de manjar, que s'apèla \VUgras{relòtge de paret}
\textsuperscript{(20)}, a una forma fòrça particulara. Es encaissat dins un cadre
de fusta que conten en mai un \VUgras{baromètre} \textsuperscript{(19)} e un
\VUgras{termomètre} \textsuperscript{(18)}.

%%K t06-tit-6 sub
\VUsaut{4.6mm}
\VUpk{10.2pt}{40}{Lo Servici del sopar.}
\VUsaut{3.0mm}

%%K t06-c3-05-1 p
5. --- Totas las parets de la sala son revestidas de \VUgras{lambris de fusta}
\textsuperscript{(21)} de garric cerat. Una \VUgras{portièra}
\textsuperscript{(22)} de tela barra pro plan la pòrta que dona sus la galariá.
Aquí anarà la familha prene lo cafè aprèp lo sopar. Las \VUgras{tassas}
\textsuperscript{(24)} e las \VUgras{sostassas} \textsuperscript{(25)} son ja
prèstas sul \VUgras{bufet de la vaissèla} \textsuperscript{(23)}.

%%K t06-c3-06-1 p
6. --- Dessús la taula i a fixat al plafon un \VUgras{lum pendolat}
\textsuperscript{(26)}, que sa lutz, fòrça viva, esclaira al vèspre tota la sala
de manjar.

%%K t06-c3-07-1 p
7. --- Ocupem-nos ara de la \VUgras{taula de manjar} \textsuperscript{(27)}
meteissa. Quand dintrèt la familha, la taula èra parada. Sus la \VUgras{toalha}
\textsuperscript{(28)} s'èran pausats, a la plaça de cada convidat, un
\VUgras{plat} \textsuperscript{(33)}, una \VUgras{servieta}
\textsuperscript{(29)}, una \VUgras{culhièra} \textsuperscript{(34)}, una
\VUgras{forceta} \textsuperscript{(35)}, un \VUgras{cotèl}
\textsuperscript{(36)} e un \VUgras{veire} \textsuperscript{(32)}. I aviá tanben
de \VUgras{botelhas} \textsuperscript{(30)} pel vin e una \VUgras{carafa}
\textsuperscript{(31)} per l'aiga.

%%K t06-c3-08-1 p
8. --- Lo \VUgras{sirvent} \textsuperscript{(39)} portèt d'en primièr la
\VUgras{sopièra} \textsuperscript{(37)}, ont fumeja encara un pauc de sopa. Ara
servís lo primièr \VUgras{plat} \textsuperscript{(38)}, un plat de carn. Puèi
presentarà los qu'avèm ja nomenats e, enfin, aprèp las \VUgras{postres},
profitarà que sos mèstres seràn passats dins la galariá per desparar la taula e
arrengar la sala de manjar.

%%K t06-c3-08-2 p
\textit{Nòta.} --- \textit{Los mots corrents relatius al manjar se donan aicí
d'un biais fòrça somari. La lista se completarà dins las doas taulas «L'Ostalariá»
(n° 13) e «Lo Mercat» (n° 15).}

%%K t06-tit-7 sub
\VUsaut{4.6mm}
\VUpk{10.2pt}{40}{La Cosina.}
\VUsaut{3.0mm}

%%K t06-c4-01-1 p
1. --- Dins la cosina, ont donam una uèlhada per la pòrta entredobèrta, trobam un
desòrdre pintoresc. Davant lo \VUgras{fogal} \textsuperscript{(1)}, ont
belugueja lo \VUgras{fuòc} \textsuperscript{(3)}, es montat lo \VUgras{ast
giratòri} \textsuperscript{(5)}. Un bèl \VUgras{rostit} \textsuperscript{(7)} es a
l'\VUgras{ast} \textsuperscript{(6)} e acaba de se far.

%%K t06-c4-02-1 p
2. --- Lo \VUgras{bofet} \textsuperscript{(8)} e la \VUgras{pala del carbon}
\textsuperscript{(9)}, que venon d'èsser emplegats, son demorats per tèrra, coma
los \VUgras{socs} \textsuperscript{(10)}.

%%K t06-c4-03-1 p
3. --- L'estagièra de la chiminèa es arrengada: lo \VUgras{fanal}
\textsuperscript{(11)}, lo \VUgras{portalumetas} \textsuperscript{(13)}, amb las
\VUgras{lumetas} \textsuperscript{(12)} dedins, lo \VUgras{candelièr}
\textsuperscript{(14)} e la \VUgras{bogeta} \textsuperscript{(15)} amb sa
\VUgras{candela} \textsuperscript{(16)} son a lor plaça. Mas lo grand \VUgras{selh
del carbon} \textsuperscript{(18)}, plen del \VUgras{carbon de pèira}
\textsuperscript{(17)} que la \VUgras{cosinièra} \textsuperscript{(23)} ven
d'anar quèrre a la cava, es demorat al mitan de la cosina.

%%K t06-c4-04-1 p
4. --- La vertat es que la paura femna se deu despachar per que lo dinnar siá
prèst a temps. Ara es contra lo \VUgras{fornèl} \textsuperscript{(19)}, a
remenar una \VUgras{sauça} \textsuperscript{(24)} que deu pas daissar cramar.

%%K t06-c4-05-1 p
5. --- Dins lo \VUgras{forn} \textsuperscript{(20)} còi un gastèl qu'a preparat
per lo grand-gost dels enfants. Dessús lo fornèl pendolan la \VUgras{culhièra
granda} \textsuperscript{(21)} e l'\VUgras{escumadoira} \textsuperscript{(22)}.

%%K t06-tit-8 sub
\VUsaut{4.0mm}
\VUpk{10.2pt}{40}{La Bateria de cosina.}
\VUsaut{3.0mm}

%%K t06-c4-06-1 p
6. --- La \VUgras{bateria de cosina} \textsuperscript{(25)}, pendolada al fons de
la sala, es fòrça neta. Las bèlas \VUgras{casseròlas} \textsuperscript{(26)} de
coire, lo \VUgras{lachièr} \textsuperscript{(27)}, lo \VUgras{colador}
\textsuperscript{(28)} e la \VUgras{raspa} \textsuperscript{(29)} son estats
netejats amb suènh.

%%K t06-c4-07-1 p
7. --- Lo \VUgras{salièr} \textsuperscript{(30)} es pausat sul pichon bufet, contra
la \VUgras{tetièra} \textsuperscript{(31)} e la \VUgras{damajana de l'òli}
\textsuperscript{(32)}. Lo \VUgras{tirador} \textsuperscript{(33)} garda
segurament los cobèrts e los cotèls de cosina.

%%K t06-c4-08-1 p
8. --- L'\VUgras{escoba} \textsuperscript{(34)} e lo \VUgras{plumalh}
\textsuperscript{(35)} son dins un canton. La cosinièra ven d'emplegar lo
\VUgras{talhador} \textsuperscript{(36)}; l'a daissat contra la fenèstra.

%%K t06-c4-09-1 p
9. --- Una \VUgras{toalha de mans} \textsuperscript{(37)} pendola a la paret,
contra l'\VUgras{embut} \textsuperscript{(38)}. Dins aquesta partida de la cosina
se gardan la \VUgras{grasilha} \textsuperscript{(39)}, lo
\VUgras{cotèl de picar} \textsuperscript{(40)} e la \VUgras{pòst de picar}
\textsuperscript{(41)}. Puèi, dessús lo cotèl, sus una \VUgras{estagièra}
\textsuperscript{(42)}, i a d'\VUgras{olas de terrilha} \textsuperscript{(43)},
l'\VUgras{ola} \textsuperscript{(44)} amb son \VUgras{cobèrcle}
\textsuperscript{(45)} e la \VUgras{marmita} \textsuperscript{(46)}. La
\VUgras{padena} \textsuperscript{(47)} pendola al costat.

%%K t06-c4-10-1 p
10. --- Sol lo \VUgras{robinet} \textsuperscript{(48)} de coire pausa una lusor
dins aqueste canton. La \VUgras{pila} \textsuperscript{(49)}, sus la quala i a lo
grand \VUgras{dorc} \textsuperscript{(50)}, deu èsser fòrça comòda amb son
armariet, ont se gardan lo \VUgras{barral del vinagre} \textsuperscript{(51)} amb
son \VUgras{canèl} \textsuperscript{(52)}, la \VUgras{caissa de las bordilhas}
\textsuperscript{(53)} e la \VUgras{vaissèla bruta} \textsuperscript{(54)}.

%%K t06-tit-9 sub
\VUsaut{4.0mm}
\VUpk{10.2pt}{40}{Autras causas que se gardan dins la cosina.}
\VUsaut{3.0mm}

%%K t06-c4-11-1 p
11. --- Lo grand \VUgras{cussòl} \textsuperscript{(55)} ont se lavan las
\VUgras{legumas} \textsuperscript{(58)} e lo \VUgras{caudièr}
\textsuperscript{(56)} de fonta pausat sul \VUgras{carrelatge}
\textsuperscript{(57)} devon tanben aver lor plaça dins aquel armari.

%%K t06-c4-12-1 p
12. --- A la cosinièra li mancan pas los fornelets, que vaicí en mai, al canton de
la taula, un \VUgras{fornelet de gas} \textsuperscript{(59)} ont s'escalfa d'aiga
dins una casseròla pichona. La \VUgras{vapor} \textsuperscript{(60)} que ne
sortís nos indica que deu bolir. Segurament aquela aiga servirà pel cafè, que
vaicí, a l'autre cap de la taula, la \VUgras{cafetièra} \textsuperscript{(64)} e
lo \VUgras{molin de cafè} \textsuperscript{(63)}.

%%K t06-c4-13-1 p
13. --- Lo fornelet es unit al \VUgras{bec de gas} \textsuperscript{(62)} per un
long \VUgras{tuble de goma} \textsuperscript{(61)}.

%%K t06-c4-14-1 p
14. --- L'\VUgras{ajudarèla} \textsuperscript{(66)}, que ven ajudar la cosinièra,
prepara las legumas pel sopar. Quand seràn pelhadas e lavadas, las pausarà dins
la \VUgras{cassòla} \textsuperscript{(65)} fins a l'ora de las còire.

%%K t06-tit-10 sub
\VUsaut{4.0mm}
{\centering\fontsize{10.2pt}{10.2pt}\selectfont\textls[40]{\scshape La Sala de Banh.} \textit{(La cambra del banh.)}\par}
\VUsaut{3.0mm}

%%K t06-c5-01-1 p
1. --- Nos demòra pas qu'una indiscrecion a far per conéisser las principalas
cambras de l'ostal: veire l'installacion de la sala de banh. Serà pas long, qu'es
pas granda, e sauprem lèu que la \VUgras{banhadoira} \textsuperscript{(1)} a un
bon \VUgras{escalfador} \textsuperscript{(2)}, que lo \VUgras{portasabon}
\textsuperscript{(3)} es a portada de man de qual se banha e que lo
\VUgras{portatoalhas} \textsuperscript{(5)} permet d'eissugar fòrça aisidament
las \VUgras{toalhas} \textsuperscript{(4)}.

%%K t06-c5-02-1 p
2. --- Aquesta sala a las paretas cobèrtas de \VUgras{cairons}
\textsuperscript{(6)}, çò qu'a permés d'installar una \VUgras{docha}
\textsuperscript{(7)} sens crénher que l'aiga qu'esposca en la fasent servir
degalhe las parets. Una pichona \VUgras{cuveta} \textsuperscript{(8)} de zinc, que
sèrv pels banhs de pès, completa lo moblament.
\VUsaut{6.0mm}
\VUfilet{22mm}
